\par\clearpage
\noindent{\small\textbf{ZUR HANNOVERSCHEN TRIANGULATION.}}\hfill359\par\medskip
\begin{center}
\textsc{Gauss} an \textsc{Bessel}.\quad Göttingen, 5.~November 1823.
\end{center}
\begingroup
\makeatletter
\renewcommand{\@makefnmark}{\hbox{\normalfont[\textsuperscript{*})]}}
\long\def\@makefntext#1{\parindent1.5em\noindent[\textsuperscript{*})\ #1}
\makeatother

\par\noindent\hspace*{\parindent}.\ .\ .\ .\ .\ .\quad Einen Teil des Jahres habe ich damit verbracht, den noch ausstehenden Teil meiner trigonometrischen Messungen im Norden von Timpenberg, Nindorf und Lüneburg bis Hamburg abzuschließen. Dann habe ich vorläufig auch die weiter westlich gelegene Gegend in Richtung Bremen erkundet, da unsere Regierung wünscht, dass die Messungen weiter nach Westen bis zur niederländischen Grenze und zum Anschluss an die \textsc{Krayenhoff}schen Dreiecke fortgesetzt werden. Schließlich habe ich noch einmal den Brocken und den Hohehagen besucht, weil zum einen die Winkel des Dreiecks Hohehagen-Hils-Brocken im Jahr 1821 unter sehr ungünstigen Umständen gemessen worden waren und zum anderen noch die damals misslungene Verbindung des Brockens mit dem Inselsberg hergestellt werden musste. Auch der hessische Dreieckspunkt Meisner sollte noch angeschlossen werden.
Ich habe diese Ziele \textls[120]{größtenteils} zu meiner Zufriedenheit erreicht. Nur den Hils hätte ich gern auch noch einmal besucht, um den Winkel dort genauer zu messen, doch die fortgeschrittene Jahreszeit hat mich daran gehindert.
Von diesem Umstand abgesehen, kann ich jetzt die Triangulation für die Gradmessung, soweit sie in meinen Zuständigkeitsbereich fällt, als abgeschlossen betrachten.
Weitere astronomische Beobachtungen als die zur Orientierung meiner ersten Dreiecksseite sind noch nicht vorgenommen worden.
Ob \textls[120]{ich} den oben erwähnten Plan zur Fortsetzung der Messungen nach Westen noch ausführe, ist übrigens noch sehr ungewiss.
Manches spricht dafür, manches, fast noch mehr, dagegen, auch abgesehen davon, dass sich vielleicht noch die Möglichkeit einer Änderung meiner äußeren Lage ergeben könnte\footnote{Es handelte sich um die Berufung nach Berlin.]}.\par

\par\noindent\hspace*{\parindent}.\ .\ .\ .\ .\ .\quad Ich habe das System meiner Hauptdreiecke in diesen Tagen sorgfältig ausgeglichen, sodass sowohl die Winkelsumme jedes einzelnen Dreiecks als auch die Seitenverhältnisse in den gekreuzten Vierecken und Fünfecken genau übereinstimmen, und zwar ohne jede Willkür, ohne Auswählen, ohne Ausschließen, alles nach den strengen Regeln der Wahrscheinlichkeitsrechnung.
Es sind insgesamt 26 Dreiecke, in denen ich alle Winkel selbst beobachtet habe.
\par\clearpage
\noindent360\hfill{\small\textbf{BRIEFWECHSEL MIT BESSEL.}}\par\medskip
\noindent Die größte Fehlersumme beträgt \(2^{\prime\prime}{,}2\) in einem Dreieck, bei dem das Anvisieren an einer Seite sehr schwierig war. Die nächstgrößte beträgt \(1^{\prime\prime}{,}8\).
Keine der 76 vorkommenden Richtungen wurde bei der Ausgleichung um eine ganze Bogensekunde geändert. Die größte Änderung beträgt \(0^{\prime\prime}{,}813\) bei der oben erwähnten Seite von Nindorf nach Hamburg.
Was ich nach meiner neuen Wahrscheinlichkeitstheorie den mittleren Fehler nenne, beträgt bei den Richtungen \(0^{\prime\prime}{,}48\).
\par\endgroup

